\section{The ambient equation and the endpoint criteria}
Fix an integer $d\geq3$. All constants in this paper may depend on $d$.
We use the Euclidean norm on vectors, the operator norm on linear maps when
estimating compositions, and the Frobenius norm in the statement of the
singularity conclusions. Put $s=\lfloor d/2\rfloor+3$. In particular,
$s>d/2+1$, and $H^s(\R^d)$ embeds into $C^{1,\alpha}_b$ for some
$\alpha\in(0,1)$ fixed throughout this section. For smooth vector fields,
$\|u\|_{H^m}^2=\sum_{|\gamma|\leq m}\|\partial^\gamma u\|_2^2$;
replacing multi-indices by ordered words changes only finite constants.

\begin{lemma}[Local theory and persistence]\label{lem:local}
Every solenoidal datum in all integer Sobolev spaces has a unique smooth
Sobolev Euler solution on an interval of length bounded below in terms of
its $H^s$ norm. The solution belongs to
$C_tH^m\cap C^1_tH^{m-1}$ on every shorter closed interval, for every
$m\geq1$. For $m\geq s$ it satisfies
\begin{equation}\label{eq:tame-energy}
 \frac{d}{dt}\|u\|_{H^m}\leq C_{d,m}\|\nabla u\|_\infty\|u\|_{H^m}.
\end{equation}
The same solution is obtained at every regularity level.
\end{lemma}
\begin{proof}
Let $J_N$ be the orthogonal Fourier projection onto $\{|\xi|\le N\}$
and solve
$\partial_tu_N=-J_N\Pi((u_N\cdot\nabla)u_N)$ with
$u_N(0)=J_Nu_0$, where $\Pi$ is the solenoidal projection.
The right-hand side is locally Lipschitz on the band-limited subspace.
Both projections commute with differentiation; $J_Nu_N=u_N$ and
$\Pi u_N=u_N$. Consequently they disappear when the differentiated
equation is tested against a derivative of $u_N$. For each differentiated
transport term the top derivative cancels against its divergence-free
transport field; the pressure cancels against a solenoidal test field.
At order one the remaining integral is bounded directly by
$C_d\|\nabla u\|_\infty\|\nabla u\|_2^2$. For $m\geq2$ the
remaining terms are products with derivative orders $j$ and $m-j+1$. The Gagliardo--Nirenberg inequalities with endpoints
$\nabla u\in L^\infty$ and $D^m u\in L^2$ give, for $1\leq j\leq m$,
\[
 \|D^j u\|_{L^{2(m-1)/(j-1)}}
 \leq C_{d,m}\|\nabla u\|_\infty^{1-(j-1)/(m-1)}
                 \|D^m u\|_2^{(j-1)/(m-1)}.
\]
The exponent is interpreted as infinity at $j=1$. The two reciprocal
Lebesgue exponents sum to $1/2$, and the two interpolation parameters
sum to one. This proves the homogeneous top-order estimate. Summing the
lower orders and retaining conservation of $L^2$ energy proves
\eqref{eq:tame-energy}. The argument uses coordinate differentiation and
integration by parts, not a three-dimensional curl identity.

At order $s$, the Sobolev embedding gives $y'\leq C_d y^2$. Thus the
cutoff solutions have a common interval, depending only on the initial
$H^s$ norm. Every higher norm is bounded there by
\eqref{eq:tame-energy}. The discarded frequencies of the quadratic term
have $L^2$ norm tending to zero uniformly on shorter intervals, since
that term is bounded in arbitrarily high Sobolev spaces. The $L^2$
difference energy estimate, followed by interpolation with the uniform
higher bounds, yields convergence in $C_tH^m$ for every fixed $m$.
The projected equation then gives convergence of time derivatives in
$C_tH^{m-1}$. It also supplies the one-sided derivatives at the time
endpoints. The same difference estimate proves uniqueness. Consequently
the solutions produced using different cutoff orders or Sobolev indices
agree on their common intervals. This proves persistence and defines a
single maximal smooth Sobolev solution.
\end{proof}

\begin{lemma}[Literal Newton pressure]\label{lem:newton}
For a smooth solenoidal $u$ in all Sobolev spaces, set
$F=\sum_{i,j}\partial_i\partial_j(u_i u_j)$ and
$N_d(x)=|x|^{2-d}/((d-2)|\mathbb S^{d-1}|)$ for $x\ne0$.
Since $|\mathbb S^{d-1}|=d\,|B_1|$, this coefficient is exactly
$[d(d-2)|B_1|]^{-1}$, as in the frozen normalization. The integral
$P=N_d*F$ is absolutely convergent at every point and smooth.
Its gradient is the $L^2$ gradient part of $-(u\cdot\nabla)u$.
Thus the projected equation gives the Euler equation with exactly this
pressure, without an undetermined function of time.
\end{lemma}
\begin{proof}
Every derivative of $F$ is both bounded and integrable: expand by the
product rule and use the $L^2$ bounds for the two factors of each product.
The kernel is locally integrable and bounded outside the unit ball, so
the Newton integral is absolutely convergent. Equivalently, with
$G_t(x)=(4\pi t)^{-d/2}e^{-|x|^2/(4t)}$,
\[
 N_d*F=\int_0^\infty G_t*F\,dt.
\]
Indeed $\int_0^\infty G_t(x)\,dt=N_d(x)$ for $x\ne0$, by the substitution
$r=|x|^2/(4t)$ and the Gamma integral. Absolute Fubini follows from
$\|G_t*|D^\gamma F|\|_\infty\leq
\min(\|D^\gamma F\|_\infty,C_dt^{-d/2}\|D^\gamma F\|_1)$.
This envelope is time integrable because $d>2$. It also permits
differentiation under the integral to every spatial order.

For each coordinate derivative, the small-time $L^2$ bound is
$\|\partial_i F\|_2$, while the large-time bound is
$\|\partial_iG_t\|_2\|F\|_1\leq C_dt^{-d/4-1/2}\|F\|_1$.
The latter is integrable at infinity for $d>2$. Hence $\nabla P\in L^2$,
and the same argument applies to all its derivatives. The heat equation
and its limits at zero and infinity give $-\Delta P=F$.
Writing $B=(u\cdot\nabla)u$, incompressibility gives
$\operatorname{div}B=F$, so $\nabla P+B$ is solenoidal.
Moreover $\nabla P$ lies in the closed $L^2$ gradient space. Indeed its
distributional curl vanishes. After Fourier transformation, its $L^2$
transform is parallel to $\xi$ for almost every $\xi\ne0$; this is
exactly the range of the orthogonal gradient multiplier
$\xi\otimes\xi/|\xi|^2$. The zero frequency is a set of measure zero.
Cutting off in frequency away from zero and then approximating the
potentials by compact smooth functions identifies this range with the
closed $L^2$ gradient space. Orthogonal
Helmholtz projection therefore yields $\nabla P=\Pi B-B$.

This normalization is also preserved under the limiting procedure in
Lemma~\ref{lem:local}. More generally, suppose
$u_n\to u$ in $C([0,S];H^m)$ for every finite $m$, and define $F_n,P_n$
by the same formulas. For each multi-index $\gamma$, the product rule,
Cauchy--Schwarz, and Sobolev embedding give
\[
 \sup_{t\le S}\bigl(\|D^\gamma(F_n-F)\|_1+
                    \|D^\gamma(F_n-F)\|_\infty+
                    \|\nabla D^\gamma(F_n-F)\|_2\bigr)\longrightarrow0.
\]
Splitting the heat integral at time one in the two estimates above yields
\begin{align*}
 \|D^\gamma(P_n-P)\|_\infty
 &\le C_d\bigl(\|D^\gamma(F_n-F)\|_\infty+
                         \|D^\gamma(F_n-F)\|_1\bigr),\\
 \|\nabla D^\gamma(P_n-P)\|_2
 &\le C_d\bigl(\|\nabla D^\gamma(F_n-F)\|_2+
                         \|D^\gamma(F_n-F)\|_1\bigr).
\end{align*}
These estimates are uniform on $[0,S]$. Thus the Newton scalars converge
with every spatial derivative uniformly, and their gradients converge
in every Sobolev space. The same estimates applied to two nearby times
give their time continuity. In particular the normalization is specified
by the convolution before and after passage to the limit.
\end{proof}

\begin{lemma}[Stability on varying intervals]\label{lem:stability}
Let $u$ be a smooth Euler solution on $[0,S]$ and let $v_n$ be smooth
solutions on $[0,S_n]$, where $0<S_n\leq S$. If
$v_n(0)\to u(0)$ in $H^s$, then
\[
 \sup_{0\leq t\leq S_n}\|v_n(t)-u(t)\|_{H^s}\longrightarrow0.
\]
No uniform higher norm of $v_n$ is assumed.
\end{lemma}
\begin{proof}
Put $w=v_n-u$. Its equation is
$w_t+(u+w)\cdot\nabla w+w\cdot\nabla u+\nabla q=0$.
Differentiation through order $s$, the transport cancellation, and the
Sobolev product estimate give
\[
 y'\leq C_d\|u\|_{H^{s+1}}y+C_d y^2,
 \qquad y=\|w\|_{H^s}.
\]
All coefficients involving $u$ are bounded on $[0,S]$ independently of
$n$. While $y\leq1$, Gronwall gives $y(t)\leq y(0)e^{CS}$.
For large $n$ this is less than $1/2$; continuity excludes a first time
when $y=1$. The estimate therefore holds on the whole actual interval
$[0,S_n]$. Sobolev embedding also gives uniform convergence of the
velocity gradients there.
\end{proof}

\begin{lemma}[Continuation and the exact vorticity norm]\label{lem:continuation}
Let $u$ be the maximal solution from Lemma~\ref{lem:local}, and suppose
its maximal lifetime $T^*$ is finite. Write
$\Omega_{ij}=\partial_i u_j-\partial_j u_i$ and
$|\Omega|=(\sum_{i<j}|\Omega_{ij}|^2)^{1/2}$. Then
\[
 \limsup_{t\uparrow T^*}\|\nabla u(t)\|_{L^\infty,\mathrm F}=\infty,
 \qquad \int_0^{T^*}\|\Omega(t)\|_\infty\,dt=\infty.
\]
\end{lemma}
\begin{proof}
If $\|\nabla u\|_\infty$ were bounded, \eqref{eq:tame-energy} would
bound the $H^s$ norm up to $T^*$. Restarting the local solution at times
approaching $T^*$ gives a common positive extension length.
The same energy estimate propagates every higher norm on a smaller
common interval, so uniqueness gives an extension in the full smooth
Sobolev class, contrary to maximality. A finite endpoint limsup would
imply such a bound, since the solution is smooth on shorter intervals.
The inequalities $\|A\|_{\rm op}\leq\|A\|_{\rm F}\leq
\sqrt d\,\|A\|_{\rm op}$ identify the required Frobenius conclusion.

For the second assertion, incompressibility gives
$\sum_i\partial_i\Omega_{ij}=\Delta u_j$. Consequently
\[
 \partial_k u_j=-\sum_i\partial_k\partial_i(-\Delta)^{-1}\Omega_{ij}.
\]
Each second derivative of the Newton kernel is a principal-value kernel
homogeneous of degree $-d$, smooth on the sphere and of spherical mean
zero, together with a constant multiple of the delta distribution.
For $0<a<1/2$, split its action into $|z|<a$, $a<|z|<2$, and a
smoothly cut off far part. In the first part subtract $\Omega(x)$.
The $C^{0,\alpha}$ bound for $\Omega$ makes that part at most
$C_da^\alpha\|u\|_{H^s}$. The annular part and local delta term are at
most $C_d\|\Omega\|_\infty(1+\log(2/a))$. In the far part use
$\Omega_{ij}=\partial_i u_j-\partial_j u_i$ and integrate by parts.
The differentiated cutoff kernel belongs to $L^2$, because it decays
as $|z|^{-d-1}$; this part is bounded by $C_d\|u\|_2$.
Taking $a=c(e+\|u\|_{H^s})^{-1/\alpha}$ gives
\begin{equation}\label{eq:log-vorticity}
 \|\nabla u\|_\infty\leq
 C_d\bigl(1+\|u(0)\|_2+
       \|\Omega\|_\infty\log(e+\|u\|_{H^s})\bigr).
\end{equation}
The finite sum over $i,j,k$ is controlled by the once-counted norm
$|\Omega|$ with a constant depending only on $d$.
Set $Z=\log(e+\|u\|_{H^s})$. By \eqref{eq:tame-energy} and
\eqref{eq:log-vorticity},
$Z'\leq C_d(1+\|u(0)\|_2)+C_d\|\Omega\|_\infty Z$.
A finite vorticity integral would bound $Z$ by Gronwall, hence permit
the same restart beyond $T^*$, a contradiction.
For smooth fields the pointwise and essential suprema agree: a strict
pointwise lower bound persists on an open ball of positive measure.
Thus these are precisely the essential-supremum norms in $H_d$.
\end{proof}

\section{Full-dimensional solenoidal geometry}
Let $E=\operatorname{span}(e_1,e_2,e_3)$ and let
$\mathcal R=I_E\oplus(-I_{E^\perp})$. The fields used below have odd
parity and satisfy $u(\mathcal R x)=\mathcal R u(x)$.
They are fields on all of $\R^d$, not extensions constant in the extra
coordinates.

\begin{lemma}[Antisymmetric potentials]\label{lem:potential}
If $A_{ij}=-A_{ji}$ is a smooth compactly supported tensor and
$v_j=\sum_i\partial_iA_{ij}$, then $\operatorname{div}v=0$.
For any trace-free matrix $L$,
\[
 A_{ij}(x)=\frac{x_i(Lx)_j-x_j(Lx)_i}{d}
 \quad\Longrightarrow\quad \sum_i\partial_i A_{ij}=(Lx)_j.
\]
Multiplying this tensor by an even compact Gevrey cutoff equal to one
near zero produces an odd compact solenoidal velocity equal to $Lx$
near zero. If $L\mathcal R=\mathcal R L$ and the cutoff is radial, the
velocity is $\mathcal R$-equivariant.
\end{lemma}
\begin{proof}
Commutation of mixed derivatives and exchange of $i,j$ make
$\sum_{i,j}\partial_j\partial_iA_{ij}$ equal to its negative.
For the displayed tensor, differentiation of the first numerator term
gives $(d+1)(Lx)_j$, and differentiation of the second gives
$(Lx)_j+x_j\operatorname{tr}L$. Their difference is $d(Lx)_j$.
The cutoff preserves antisymmetry and compact support. The tensor is
even, so its divergence is odd. Under $\mathcal R$ the tensor transforms
as a two-tensor; applying divergence gives the stated equivariance.
Gevrey bounds follow from the finite product rule with a polynomial.
\end{proof}

The same replacement removes the use of curl from oscillatory
solenoidalization. Put $D_i=\kappa\partial_{y_i}+m_i\partial_\theta$.
These constant-coefficient operators commute. Thus
$\sum_iD_i\sum_jD_j A_{ji}=0$ for every antisymmetric $A$.
When $m\cdot v=0$, $|m|=1$ and $F'=f$,
\[
 A_{ij}=(m_i v_j-m_j v_i)F(\theta)
 \quad\Longrightarrow\quad
 \sum_iD_i A_{ij}=v_jf(\theta)+\kappa\sum_i\partial_{y_i}A_{ij}.
\]
The last term is the actual spatial corrector, retained also at the
terminal truncation order. This identity treats every ambient component.

\begin{lemma}[Symmetry of the ambient pressure and central blocks]\label{lem:reflection}
If $u$ is $\mathcal R$-equivariant, its Newton pressure obeys
$P(\mathcal R x)=P(x)$. At $x\in E$, both $\nabla u(x)$ and
$\nabla^2P(x)$ preserve $E$ and $E^\perp$.
\end{lemma}
\begin{proof}
The chain rule and orthogonality show that
$\sum_{i,j}\partial_i\partial_j(u_i u_j)$ is invariant under
$\mathcal R$. The Newton kernel is radial and Lebesgue measure is
orthogonally invariant, so a change of variables in its absolutely
convergent convolution proves the pressure identity. Differentiate the
velocity identity once and the pressure identity twice at a fixed point.
The resulting maps commute with $\mathcal R$, and therefore preserve its
two eigenspaces. No assertion that the active block is trace free follows
or is used: the full trace alone vanishes.
\end{proof}

\begin{lemma}[The endpoint operator preserves the active plane]\label{lem:endpoint-symmetry}
Along the central trajectory of an equivariant parent, suppose the normal
$m(t)$ belongs to $E$. By Lemma~\ref{lem:reflection}, its gradient $M(t)$
and pressure Hessian $H(t)$ commute with $\mathcal R$. For the
stationary displacement problem on $[0,t_0]$ with
$m(t)\cdot\eta(t)=0$, $\eta(0)=0$, prescribed $\eta(t_0)=Y$, and
coercive action
$\mathcal A(\eta)=\int_0^{t_0}(|\eta'|^2-\langle H\eta,\eta\rangle)dt$,
the endpoint map $\Lambda Y=\operatorname{proj}_{m(t_0)^\perp}\eta'(t_0)$
commutes with $\mathcal R$. In particular a prescribed endpoint in
$E\cap m(t_0)^\perp$ produces a displacement and velocity in $E$.
\end{lemma}
\begin{proof}
Reflection preserves the constraint, the two endpoints and the action.
Uniqueness for the coercive variational problem therefore gives
$\eta_{\mathcal RY}=\mathcal R\eta_Y$. Taking the terminal derivative
and its orthogonal projection proves the claim about $\Lambda$.
If $Y=\mathcal RY$, uniqueness also gives
$\eta_Y=\mathcal R\eta_Y$ at every time. The identity remains true for
the physical velocity $\eta'-M\eta$ because $M$ commutes with the same
reflection. Hence its extra components vanish exactly. This argument
does not infer vanishing from an estimate that could be lost under the
subsequent endpoint normalization.
\end{proof}
